\noindent\textbf{Elección del procedimiento.} Con cilindro interior fijo y exterior móvil, $u_\theta=Ar+B/r$. Las condiciones $AR_1+B/R_1=0$ y $AR_2+B/R_2=\Omega_2R_2$ dan $A=\Omega_2R_2^2/(R_2^2-R_1^2)$ y $B=-AR_1^2$. El par resistente tiene módulo $4\pi\mu L|B|$.

\begin{enumerate}

\item Eliminar $B$ con la condición interior antes de despejar $A$.

\[\fitmath{\mathrm{Acoef}=\frac{W_{2}\cdot {R_{2}}^{2}}{{R_{2}}^{2}-{R_{1}}^{2}}}\]

\[\fitmath{\mathrm{Acoef}=\frac{(20)\cdot {(0.05)}^{2}}{{(0.05)}^{2}-{(0.04)}^{2}}\approx 55.55556\,\mathrm{1/s}}\]

\item El signo garantiza velocidad nula en la pared interior.

\[\fitmath{\mathrm{Bcoef}=-\mathrm{Acoef}\cdot {R_{1}}^{2}}\]

\[\fitmath{\mathrm{Bcoef}=-(55.55555556)\cdot {(0.04)}^{2}\approx -0.08888889\,\mathrm{m^{2}/s}}\]

\item Integrar esfuerzo tangencial por brazo sobre la superficie.

\[\fitmath{T=4\cdot \pi\cdot \mu\cdot L\cdot -\mathrm{Bcoef}}\]

\[\fitmath{T=4\cdot \pi\cdot (1.41)\cdot (0.2)\cdot -(-0.08888888889)\approx 0.314997\,\mathrm{N\,m}}\]

\item El motor exterior aporta potencia que se disipa en el fluido.

\[\fitmath{P=T\cdot W_{2}}\]

\[\fitmath{P=(0.3149970234)\cdot (20)\approx 6.29994\,\mathrm{W}}\]

\end{enumerate}
